
          %%%%% ~~~~~~~~~~~~~~~~~~~~ %%%%%

\chapter{The representations of $\sl(3,\C)$}
\setcounter{theorem}{0}
\setcounter{equation}{0}

%\cite{vanDerVaart1996}
%\cite{Kosorok2008}

%\renewcommand{\theenumi}{\alph{enumi}}
%\renewcommand{\labelenumi}{\textnormal{(\theenumi)}$\;\;$}
\renewcommand{\theenumi}{\roman{enumi}}
\renewcommand{\labelenumi}{\textnormal{(\theenumi)}$\;\;$}

          %%%%% ~~~~~~~~~~~~~~~~~~~~ %%%%%

\section{$\sl(n,\C) \,\cong\, \su(n)_{\C} \,=\, \su(n) \oplus \i\cdot\su(n)$}

          %%%%% ~~~~~~~~~~~~~~~~~~~~ %%%%%

\begin{lemma}\label{glnCunC}
\mbox{}
\vskip 0.05cm
\noindent
\begin{enumerate}
\item
    $
    \left(\,\overset{{\color{white}.}}{\mathfrak{u}(n)}\,\right)
    \bigcap
    \left(\,\overset{{\color{white}.}}{\i\cdot\mathfrak{u}(n)}\,\right)
    \, = \,
    \{\,0\,\}
    $\,
\item
    $\gl(n,\C) \,\cong\, \mathfrak{u}(n)_{\C} \,=\, \mathfrak{u}(n) \oplus \i\cdot\mathfrak{u}(n)$\,
    as complex Lie algebras
\end{enumerate}
\end{lemma}
\proof
\vskip 0.1cm
\noindent
First, recall:
\begin{itemize}
\item
    $\gl(n,\C) \; = \; \C^{n \times n}$
\item
    $
    \mathfrak{u}(n)
    \; = \;
        \left\{\;
            X \in \C^{n \times n}
            \,\left\vert\;\,
            \overset{{\color{white}.}}{X^{*} \,=\, -X}
            \right.
            \,\right\}
    $
    \quad
    (see, for example, Proposition 3.24, p.58, \cite{Hall2015})
\end{itemize}
\begin{enumerate}
\item
    Note that
    \begin{equation*}
    Y \in \mathfrak{u}(n)
    \;\;\Longrightarrow\;\;
        Y^{*} \, = \, -Y
    \;\;\Longrightarrow\;\;
        \left(\,\overset{{\color{white}.}}{\i \cdot Y}\right)^{\!*}
        \; = \;
        (-\i) \cdot Y^{*}
        \; = \;
        (-\i) \cdot (-Y)
        \; = \;
        \i \cdot Y
    \end{equation*}
    Thus,
    \,$X \in \i \cdot \mathfrak{u}(n) \;\Longrightarrow\; X^{*} = X$.\,
    Hence, we see that 
    \begin{equation*}
    X \in
        \left(\,\overset{{\color{white}.}}{\mathfrak{u}(n)}\,\right)
        \bigcap
        \left(\,\overset{{\color{white}.}}{\i\cdot\mathfrak{u}(n)}\,\right)
    \;\;\Longrightarrow\;\;
        \left\{\begin{array}{c}
            X^{*} = -X
            \\
            X^{*} = {\color{white}-}X
            \end{array}\right.
    \;\;\Longrightarrow\;\;
        X = X^{*} = -X
    \;\;\Longrightarrow\;\;
        X = 0
    \end{equation*}
    This proves
    $
    \left(\,\overset{{\color{white}.}}{\mathfrak{u}(n)}\,\right)
    \bigcap
    \left(\,\overset{{\color{white}.}}{\i\cdot\mathfrak{u}(n)}\,\right)
    \, = \,
    \{\,0\,\}
    $,\,
    as required.
\item
    \textbf{Claim 1:}\;
    \,$\gl(n,\C) \;\subset\; \mathfrak{u}(n) \,\oplus\, \i\cdot\mathfrak{u}(n)$\,
    \vskip 0.1cm
    \noindent
    Proof of Claim 1:\;
    For each \,$X \in \gl(n,\C) = \C^{n \times n}$,\, we may write:
    \begin{equation*}
    X
    \; = \;
        X_{1} \,+\, \i \cdot X_{2}
    \; = \;
        \left(\dfrac{X - X^{*}}{2}\right)
        \,+\,
        \i\,\cdot\left(\dfrac{X + X^{*}}{\i \cdot 2}\right),
    \quad
    \textnormal{where}\;\;
    X_{1}
    \; := \;
        \left(\dfrac{X - X^{*}}{2}\right),
    \;\;
    X_{2}
    \; := \;
        \left(\dfrac{X + X^{*}}{\i \cdot 2}\right).
    \end{equation*}
    Next, note that \,$X_{1},\, X_{2} \in \mathfrak{u}(n)$.\,
    Indeed,
    \begin{eqnarray*}
    X_{1}^{*}
    & = &
        \left(\dfrac{X - X^{*}}{2}\right)^{\!*}
    \;\; = \;\;
        \left(\dfrac{X^{*} - X^{**}}{2}\right)
    \;\; = \;\;
        \left(\dfrac{X^{*} - X}{2}\right)
    \;\; = \;\;
        -\left(\dfrac{X - X^{*}}{2}\right)
    \;\; = \;\;
        -\,X_{1}
    \\
    X_{2}^{*}
    & = &
        \left(\dfrac{X + X^{*}}{\i \cdot 2}\right)^{\!*}
    \;\; = \;\;
        \left(\dfrac{X^{*} + X^{**}}{-\,\i \cdot 2}\right)
    \;\; = \;\;
        \left(\dfrac{X^{*} + X}{-\,\i \cdot 2}\right)
    \;\; = \;\;
        -\left(\dfrac{X + X^{*}}{\i \cdot 2}\right)
    \;\; = \;\;
        -\,X_{2}
    \end{eqnarray*}
    This shows indeed that
    \,$X_{1},\, X_{2} \in \mathfrak{u}(n)$,\,
    and completes the proof of Claim 1.
    \vskip 0.5cm
    \noindent
    We can now finish the proof: Conclusion (i) and Claim 1 together imply
    \begin{equation*}
    \gl(n,\C)
    \;\;\subset\;\;
        \mathfrak{u}(n) \,\oplus\, \i\cdot\mathfrak{u}(n)
    \;\; =: \;\;
        \mathfrak{u}(n)_{\C}
    \;\;\subset\;\;
        \C^{n \times n}
    \;\; = \;\;
        \gl(n,\C)\,,
    \end{equation*}
    from which it immediately follows that
    \,$\mathfrak{u}(n)_{\C} \,:=\, \mathfrak{u}(n) \,\oplus\, \i\cdot\mathfrak{u}(n) \,=\, \gl(n,\C)$.\,
    This completes the proof of Conclusion (ii), as well as that of the Lemma.
    \qed
\end{enumerate}

          %%%%% ~~~~~~~~~~~~~~~~~~~~ %%%%%

\vskip 0.5cm
\begin{proposition}
\mbox{}
\vskip 0.05cm
\noindent
$\sl(n,\C)$\, is isomorphic, as a complex Lie algebra,
to the complexification \,$\su(n)_{\C}$\, of \,$\su(n)$,\,
i.e., for each \,$n \in \N$,\, we have: 
\begin{equation*}
\sl(n,\C) \;\cong\; \su(n)_{\C}\,,
\quad
\textnormal{as complex Lie algebras}
\end{equation*}
\end{proposition}
\proof
Recall:
\begin{itemize}
\item
    $
    \sl(n,\C)
    \; = \;
        \left\{\;
            X \in \C^{n \times n}
            \,\left\vert\;\,
            \overset{{\color{white}.}}{\trace(X) = 0}
            \right.
            \,\right\}
    $
    \quad
    (see, for example, Example 3.4, p.50, \cite{Hall2015})
\item
    $
    \mathfrak{u}(n)
    \; = \;
        \left\{\;
            X \in \C^{n \times n}
            \,\left\vert\;\,
            \overset{{\color{white}.}}{X^{*} \,=\, -X}
            \right.
            \,\right\}
    $
    \quad
    (see, for example, Proposition 3.24, p.58, \cite{Hall2015})
\item
    $
    \su(n)
    \; = \;
        \left\{\;
            X \in \C^{n \times n}
            \,\left\vert
            \begin{array}{c}
                X^{*} \,=\, -X
                \\
                \trace(X) \,= 0
                \end{array}
            \right.
            \!\!\right\}
    $
    \quad
    (see, for example, Proposition 3.24, p.58, \cite{Hall2015})
\end{itemize}
Hence,
\;$
\su(n) \,=\, \mathfrak{u}(n) \,\bigcap\, \sl(n,\C)
$.\;
\vskip 0.5cm
\noindent
\textbf{Claim 1:}\;
$
\left(\,\overset{{\color{white}.}}{\su(n)}\,\right)
\bigcap
\left(\,\overset{{\color{white}.}}{\i\cdot\su(n)}\,\right)
\, = \,
    \{\,0\,\}
$\,
\vskip 0.1cm
\noindent
Proof of Claim 1:\;
First note that
\begin{equation*}
Y \in \su(n)
\;\;\Longrightarrow\;\;
    Y^{*} \, = \, -Y
\;\;\Longrightarrow\;\;
    \left(\,\overset{{\color{white}.}}{\i \cdot Y}\right)^{\!*}
    \; = \;
    (-\i) \cdot Y^{*}
    \; = \;
    (-\i) \cdot (-Y)
    \; = \;
    \i \cdot Y
\end{equation*}
Thus,
\,$X \in \i \cdot \su(n) \;\Longrightarrow\; X^{*} = X$.\,
Hence, we see that 
\begin{equation*}
X \in
    \left(\,\overset{{\color{white}.}}{\su(n)}\,\right)
    \bigcap
    \left(\,\overset{{\color{white}.}}{\i\cdot\su(n)}\,\right)
\;\;\Longrightarrow\;\;
    \left\{\begin{array}{c}
        X^{*} = -X
        \\
        X^{*} = {\color{white}-}X
        \end{array}\right.
\;\;\Longrightarrow\;\;
    X = X^{*} = -X
\;\;\Longrightarrow\;\;
    X = 0
\end{equation*}
This proves Claim 1.
\vskip 0.5cm
\noindent
\textbf{Claim 2:}\; $\su(n) \oplus \i\cdot\su(n) \,\subset\, \sl(n,\C)$\,
\vskip 0.1cm
\noindent
Proof of Claim 2:\;
Let \,$X,\, Y \in \su(n)$\, and \,$Z = X + \i\cdot Y \in \su(n) \oplus \i\cdot\su(n)$.\,
Then, \,$\trace(Z) \,=\, \trace(X) + \i\cdot\trace(Y) = 0 + \i\cdot 0 = 0$,\,
which implies that \,$Z \in \sl(n,\C)$.\,
This proves Claim 2.
\vskip 0.5cm
\noindent
\textbf{Claim 3:}\; $\sl(n,\C) \,\subset\, \su(n) \oplus \i\cdot\su(n)$\,
\vskip 0.1cm
\noindent
Proof of Claim 3:\;
Let \,$X \in \sl(n,\C)$.\,
Since \,$\sl(n,\C) \subset \gl(n,\C) = \mathfrak{u}(n) \oplus \i\cdot\mathfrak{u}(n)$,\,
by Lemma \ref{glnCunC}, we may write:
\begin{equation*}
X
\; = \;
    X_{1} \,+\, \i \cdot X_{2}\,,
\quad
\textnormal{where}\;\;
X_{1}
\, \in \,
    \mathfrak{u}(n),
\;\;
X_{2}
\, \in \,
    \mathfrak{u}(n)
\end{equation*}
Hence, Claim 3 will follow once we establish that
\,$X_{1},\, X_{2} \,\in\, \su(n)$.\,
Since we already have
\,$X_{1},\, X_{2} \,\in\, \mathfrak{u}(n)$,\,
it remains only to prove that
\,$\trace(X_{1}) \,=\, \trace(X_{2}) \,=\, 0$.\,
To this end, first note that
\,$X_{1} \in \mathfrak{u}(n) \;\Longrightarrow\; \trace(\,X_{1}\,) \in \i\cdot\Re$.\,
Indeed,
\begin{eqnarray*}
X_{1} \in \mathfrak{u}(n)
& \Longleftrightarrow &
    X_{1}^{*} \,=\, -X_{1}
\\
& \Longrightarrow &
    \trace(\,X_{1}^{*}\,) \,=\, \trace(\,-X_{1}\,)
\;\;\Longrightarrow\;\;
    \overline{\trace(\,X_{1}\,)} \,=\, -\,\trace(\,X_{1}\,)
\\
& \Longrightarrow &
    \trace(\,X_{1}\,) \in \i\cdot\Re
\end{eqnarray*}
Similarly, we have \,$\trace(\,X_{2}\,) \in \i\cdot\Re$.\,
Now, recall that \,$X_{1} + \i \cdot X_{2} = X \in \sl(2,\C)$;\,
in particular, \,$\trace(X) \,=\, 0$.\,
Hence,
\begin{eqnarray*}
0
& = &
    \trace(\,X\,)
\;\; = \;\;
    \trace(\,X_{1} + \i \cdot X_{2}\,)
\;\; = \;\;
    \trace(\,X_{1}\,) + \i\cdot\trace(\,X_{2}\,)\,,
\end{eqnarray*}
which implies
\,$
\trace(\,X_{1}\,)
\, = \,
    -\,\i\cdot\trace(\,X_{2}\,)
\, \in \,
    \left(\,\i\cdot\Re\,\right)
    \bigcap
    \left(\,\Re\,\right)
\, = \,
    \{\,0\,\}
$,\,
which in turn implies
\,$\trace(\,X_{1}\,) = \trace(\,X_{2}\,) \,=\, 0$.\,
This completes the proof of Claim 3.
\vskip 0.5cm
\noindent
The Proposition now follows immediately from Claim 2 and Claim 3.
\qed

          %%%%% ~~~~~~~~~~~~~~~~~~~~ %%%%%

\section{A common basis over \,$\C$\, for \,$\sl(3,\C)$\, and the commutation relations of the basis elements}

          %%%%% ~~~~~~~~~~~~~~~~~~~~ %%%%%

\begin{lemma}
\mbox{}
\vskip 0.05cm
\noindent
The following matrices form a basis over \,$\C$\, for 
\,$\sl(3,\C)$:
\begin{equation*}
\begin{array}{ccccccccc}
H_{1}
\!\!\!
&=&
    \!\!
    {\footnotesize\left(\,
        \begin{array}{rrr} 1 & 0 & {\color{white}..}0 \\ 0 & -1 & 0 \\ 0 & 0 & 0 \end{array}
        \,\right)},
&
H_{2}
\!\!\!
&=&
    \!\!
    {\footnotesize\left(\,
        \begin{array}{rrr} 0 & {\color{white}..}0 & 0 \\ 0 & 1 & 0 \\ 0 & 0 & -1 \end{array}
        \,\right)},
\\ \\
X_{1}
\!\!\!
&=&
    \!\!
    {\footnotesize\left(\,
        \begin{array}{rrr} 0 & {\color{white}..}1 & {\color{white}..}0 \\ 0 & 0 & 0 \\ 0 & 0 & 0 \end{array}
        \,\right)},
&
X_{2}
\!\!\!
&=&
    \!\!
    {\footnotesize\left(\,
        \begin{array}{rrr} 0 & {\color{white}..}0 & {\color{white}..}0 \\ 0 & 0 & 1 \\ 0 & 0 & 0 \end{array}
        \,\right)},
&
X_{3}
\!\!\!
&=&
    \!\!
    {\footnotesize\left(\,
        \begin{array}{rrr} 0 & {\color{white}..}0 & {\color{white}..}1 \\ 0 & 0 & 0 \\ 0 & 0 & 0 \end{array}
        \,\right)},
\\ \\
Y_{1}
\!\!\!
&=&
    \!\!
    {\footnotesize\left(\,
        \begin{array}{rrr} 0 & {\color{white}..}0 & {\color{white}..}0 \\ 1 & 0 & 0 \\ 0 & 0 & 0 \end{array}
        \,\right)},
&
Y_{2}
\!\!\!
&=&
    \!\!
    {\footnotesize\left(\,
        \begin{array}{rrr} 0 & {\color{white}..}0 & {\color{white}..}0 \\ 0 & 0 & 0 \\ 0 & 1 & 0 \end{array}
        \,\right)},
&
Y_{3}
\!\!\!
&=&
    \!\!
    {\footnotesize\left(\,
        \begin{array}{rrr} 0 & {\color{white}..}0 & {\color{white}..}0 \\ 0 & 0 & 0 \\ 1 & 0 & 0 \end{array}
        \,\right)}.
\end{array}
\end{equation*}
Furthermore,
\,$H_{1},H_{2},\; X_{1},X_{2},X_{3},\; Y_{1},Y_{2},Y_{3} \,\in\, \sl(3,\C)$\,
satisfy the following commutation relations:
\begin{equation*}
\begin{array}{ccr}
\left[H_{1},X_{1}\right] &=& 2\,X_{1},
\\
\left[H_{1},\;Y_{1}\right] &=& -\,2\,\,Y_{1},
\\
\left[X_{1},\;Y_{1}\right] &=& H_{1},
\end{array}
\quad\quad
\begin{array}{ccr}
\left[H_{2},X_{2}\right] &=& 2\,X_{2},
\\
\left[H_{2},\;Y_{2}\right] &=& -\,2\,\,Y_{2},
\\
\left[X_{2},\;Y_{2}\right] &=& H_{2},
\end{array}
\end{equation*}
\begin{equation*}
\begin{array}{ccr}
\left[H_{1},H_{2}\right] &=& 0,
\\ \\
\left[H_{1},X_{1}\right] &=& 2\,X_{1},
\\
\left[H_{2},X_{1}\right] &=& -\,X_{1},
\\ \\
\left[H_{1},X_{2}\right] &=& -\,X_{2},
\\
\left[H_{2},X_{2}\right] &=& 2\,X_{2},
\\ \\
\left[H_{1},X_{3}\right] &=& X_{3},
\\
\left[H_{2},X_{3}\right] &=& X_{3},
\end{array}
\quad\quad
\begin{array}{ccr}
&&
\\ \\
\left[H_{1},Y_{1}\right] &=& -\,2\,Y_{1},
\\
\left[H_{2},Y_{1}\right] &=& Y_{1},
\\ \\
\left[H_{1},Y_{2}\right] &=& Y_{2},
\\
\left[H_{2},Y_{2}\right] &=& -\,2\,Y_{2},
\\ \\
\left[H_{1},Y_{3}\right] &=& -\,Y_{3},
\\
\left[H_{2},Y_{3}\right] &=& -\,Y_{3},
\end{array}
\end{equation*}
\begin{equation*}
\begin{array}{ccr}
\left[X_{1},Y_{1}\right] &=& H_{1},
\\
\left[X_{2},Y_{2}\right] &=& H_{2},
\\
\left[X_{3},Y_{3}\right] &=& H_{1}+H_{2},
\\ \\
\left[X_{1},X_{2}\right] &=& X_{3},
\\
\left[X_{1},Y_{2}\right] &=& 0,
\\ \\
\left[X_{1},X_{3}\right] &=& 0,
\\
\left[X_{2},X_{3}\right] &=& 0,
\\ \\
\left[X_{2},Y_{3}\right] &=& Y_{1},
\\
\left[X_{1},Y_{3}\right] &=& -\,Y_{2},
\end{array}
\quad\quad
\begin{array}{ccr}
&&
\\
&&
\\
&&
\\ \\
\left[Y_{1},Y_{2}\right] &=& -\,Y_{3},
\\
\left[X_{2},Y_{1}\right] &=& 0,
\\ \\
\left[Y_{1},Y_{3}\right] &=& 0,
\\
\left[Y_{2},Y_{3}\right] &=& 0,
\\ \\
\left[X_{3},Y_{2}\right] &=& X_{1},
\\
\left[X_{3},Y_{1}\right] &=& -\,X_{2},
\end{array}
\end{equation*}
\end{lemma}
\proof
Recall that:
\,$
\sl(n,\C)
\; = \;
    \left\{\;
        X \in \C^{n \times n}
        \,\left\vert\;\,
        \overset{{\color{white}.}}{\trace(X) = 0}
        \right.
        \,\right\}
$\,
(see, for example, Example 3.4, p.50, \cite{Hall2015}).
Hence, \,$\dim_{\C}\!\left(\,\overset{{\color{white}.}}{\sl(3,\C)}\,\right) \,=\, 9-1 \,=\, 8$.\,
On the other hand, it is clear that
\,$H_{1}, H_{2},\; X_{1}, X_{2}, X_{3},\; Y_{1}, Y_{2}, Y_{3} \,\in\, \sl(3,\C)$,\,
and that these $8$ matrices are linearly independent.
It thus follows that
\,$H_{1}, H_{2},\; X_{1}, X_{2}, X_{3},\; Y_{1}, Y_{2}, Y_{3} \,\in\, \sl(3,\C)$\,
indeed form a basis over \,$\C$\, for \,$\sl(3,\C)$.\,
Their commutation relations follow by straightforward computations.
\qed

          %%%%% ~~~~~~~~~~~~~~~~~~~~ %%%%%

          %%%%% ~~~~~~~~~~~~~~~~~~~~ %%%%%

